% ============================================================
% CHAPTER 4: MULTIWAY CUT AND k-CUT (matches Vazirani Ch. 4)
% ============================================================
\setcounter{chapter}{3}
\chapter{Multiway Cut and k-Cut}
\label{ch:multiway}

\section{Intuition: Isolating Cuts}

\begin{intuitionbox}
\textbf{Multiway Cut:} We have $k$ terminals that must be separated. For each terminal $s_i$, find the minimum cut isolating $s_i$ from all other terminals (merge other terminals into a super-terminal). Take the union of the $k-1$ lightest isolating cuts.

\textbf{Why it works:} In the optimal multiway cut, each terminal sits in its own component. The cut separating component $i$ from the rest is an isolating cut for $s_i$. The sum of all $k$ isolating cuts in OPT is exactly $2 \times$ OPT (each edge counted twice). Discarding the heaviest leaves at most $2(1-1/k)\OPT$.
\end{intuitionbox}

\section{Multiway Cut}

\begin{definition}[Multiway Cut]
Given graph $G = (V, E)$ with edge weights and terminals $S = \{s_1, \dots, s_k\}$, find minimum-weight edge set whose removal disconnects all terminals from each other.
\end{definition}

\begin{algorithmdef}[Multiway Cut — Factor $2-2/k$]
\begin{enumerate}
    \item For each $i = 1..k$, compute minimum isolating cut $C_i$ separating $s_i$ from $S \setminus \{s_i\}$
    \item Discard the heaviest $C_i$
    \item Output union of remaining $k-1$ cuts
\end{enumerate}
\end{algorithmdef}

\begin{theorem}
This algorithm achieves a $2 - \frac{2}{k}$ approximation for Multiway Cut.
\end{theorem}

\subsection{Minimum k-Cut}

\begin{definition}[$k$-Cut]
Find minimum-weight edge set whose removal partitions $G$ into at least $k$ components.
\end{definition}

The Gomory-Hu tree represents all pairwise minimum cuts in a tree with $n-1$ edges. Removing the $k-1$ lightest edges gives a $k$-cut.

\begin{algorithmdef}[Minimum $k$-Cut — Factor $2-2/k$]
\begin{enumerate}
    \item Compute Gomory-Hu tree $T$ of $G$
    \item Remove $k-1$ lightest edges from $T$
    \item Output union of corresponding cuts in $G$
\end{enumerate}
\end{algorithmdef}

\begin{theorem}
This achieves a $2 - \frac{2}{k}$ approximation for Minimum $k$-Cut.
\end{theorem}

\subsection{Other Approaches}

\begin{table}[H]
\centering
\caption{Multiway Cut / k-Cut Algorithms}
\begin{tabular}{lccc}
\toprule
Algorithm & Approximation & Technique \\
\midrule
Isolating cuts (this chapter) & $2-2/k$ & Min-cut subroutine \\
Gomory-Hu tree (this chapter) & $2-2/k$ & Min-cut + tree \\
LP rounding (Ch. 19) & $1.5$ & LP relaxation \\
Region growing & $2$ & Continuous relaxation \\
\bottomrule
\end{tabular}
\end{table}

\section{Python Implementation}

\begin{lstlisting}[style=python, caption=Multiway Cut and k-Cut]
import heapq


class MaxFlow:
    """Edmonds-Karp max flow / min cut on an undirected graph."""
    def __init__(self, n):
        self.n = n
        self.adj = [[] for _ in range(n)]
        self.cap = [[0] * n for _ in range(n)]

    def add_edge(self, u, v, c):
        self.cap[u][v] += c
        if v not in self.adj[u]:
            self.adj[u].append(v)
        if u not in self.adj[v]:
            self.adj[v].append(u)

    def bfs(self, s, t, parent):
        parent[:] = [-1] * self.n
        parent[s] = -2
        q = [(s, float('inf'))]
        while q:
            u, flow = q.pop(0)
            for v in self.adj[u]:
                if parent[v] == -1 and self.cap[u][v] > 0:
                    parent[v] = u
                    nf = min(flow, self.cap[u][v])
                    if v == t:
                        return nf
                    q.append((v, nf))
        return 0

    def max_flow(self, s, t):
        flow = 0
        parent = [-1] * self.n
        while True:
            nf = self.bfs(s, t, parent)
            if nf == 0:
                break
            flow += nf
            v = t
            while v != s:
                u = parent[v]
                self.cap[u][v] -= nf
                self.cap[v][u] += nf
                v = u
        return flow

    def reachable(self, s):
        seen, stack = {s}, [s]
        while stack:
            u = stack.pop()
            for v in self.adj[u]:
                if self.cap[u][v] > 0 and v not in seen:
                    seen.add(v)
                    stack.append(v)
        return seen


def min_s_t_cut(graph, s, t):
    """Min s-t cut; returns (S, T, weight)."""
    vertices = list(graph.keys())
    idx = {v: i for i, v in enumerate(vertices)}
    mf = MaxFlow(len(vertices))
    for u in graph:
        for v, c in graph[u].items():
            if u < v:
                mf.add_edge(idx[u], idx[v], c)
                mf.add_edge(idx[v], idx[u], c)
    mf.max_flow(idx[s], idx[t])
    S = {vertices[i] for i in mf.reachable(idx[s])}
    T = set(vertices) - S
    w = sum(graph[u][v] for u in S for v in T if v in graph[u])
    return S, T, w


def isolating_cut(graph, s, other_terminals):
    """Min cut separating s from all other terminals (super-source trick)."""
    vertices = list(graph.keys())
    idx = {v: i for i, v in enumerate(vertices)}
    n = len(vertices)
    INF = float('inf')
    mf = MaxFlow(n + 1)
    for u in graph:
        for v, c in graph[u].items():
            if u < v:
                mf.add_edge(idx[u], idx[v], c)
                mf.add_edge(idx[v], idx[u], c)
    for t in other_terminals:
        mf.add_edge(n, idx[t], INF)          # super-source -> other terminals
    mf.max_flow(n, idx[s])
    S = {vertices[i] for i in mf.reachable(n) if i != n}
    T = set(vertices) - S
    edges = {(min(u, v), max(u, v)) for u in S for v in T if v in graph[u]}
    w = sum(graph[u][v] for u in S for v in T if v in graph[u])
    return S, T, edges, w


def multiway_cut_2_2k(graph, terminals):
    """2-2/k approx for Multiway Cut."""
    k = len(terminals)
    if k <= 2:
        s, t = list(terminals)
        S, T, w = min_s_t_cut(graph, s, t)
        edges = {(min(u, v), max(u, v)) for u in S for v in T if v in graph[u]}
        return edges, w

    terminals = list(terminals)
    cuts, weights = [], []
    for i, s in enumerate(terminals):
        other = [t for j, t in enumerate(terminals) if j != i]
        _, _, edges, w = isolating_cut(graph, s, other)
        cuts.append(edges)
        weights.append(w)

    skip = max(range(k), key=lambda i: weights[i])
    result_edges = set()
    result_weight = 0.0
    for i in range(k):
        if i != skip:
            result_edges.update(cuts[i])
            result_weight += weights[i]
    return result_edges, result_weight


def gomory_hu_tree(graph):
    """Gomory-Hu tree; returns (edges, weights) keyed by (u, v)."""
    vertices = list(graph.keys())
    n = len(vertices)
    if n <= 1:
        return [], {}
    idx = {v: i for i, v in enumerate(vertices)}
    parent = {i: 0 for i in range(1, n)}
    tree_edges, tree_weights = [], {}
    for i in range(1, n):
        s, t = i, parent[i]
        mf = MaxFlow(n)
        for u in graph:
            for v, c in graph[u].items():
                if u < v:
                    mf.add_edge(idx[u], idx[v], c)
                    mf.add_edge(idx[v], idx[u], c)
        mf.max_flow(s, t)
        visited = mf.reachable(s)
        u, v = vertices[s], vertices[t]
        w = sum(graph[vertices[x]][vertices[y]]
                for x in visited for y in range(n)
                if y not in visited and vertices[y] in graph[vertices[x]])
        tree_edges.append((u, v))
        tree_weights[(min(u, v), max(u, v))] = w
        for j in range(i + 1, n):
            if parent[j] == t and j in visited:
                parent[j] = i
    return tree_edges, tree_weights


def min_k_cut_2_2k(graph, k):
    """2 - 2/k approx for Minimum k-Cut."""
    if k <= 1:
        return set(), 0.0
    if k >= len(graph):
        edges = {(min(u, v), max(u, v)) for u in graph for v in graph[u] if u < v}
        w = sum(graph[u][v] for u in graph for v in graph[u] if u < v)
        return edges, w

    tree_edges, tree_weights = gomory_hu_tree(graph)
    key = lambda e: tree_weights[(min(e[0], e[1]), max(e[0], e[1]))]
    to_remove = sorted(tree_edges, key=key)[:k - 1]

    result_edges, result_weight = set(), 0.0
    for u, v in to_remove:
        S, T, w = min_s_t_cut(graph, u, v)
        for x in S:
            for y in T:
                if y in graph[x]:
                    result_edges.add((min(x, y), max(x, y)))
        result_weight += w
    return result_edges, result_weight
\end{lstlisting}

\section{Applications}
\begin{itemize}
    \item \textbf{Image segmentation:} Multiway cut = partition pixels into regions
    \item \textbf{Network reliability:} $k$-cut = minimum edges to disconnect into $k$ parts
    \item \textbf{Clustering:} Multiway cut with similarity weights
    \item \textbf{VLSI partitioning:} Divide circuit into modules
\end{itemize}

\subsection*{Real Example: Medical Image Segmentation}
\textbf{Scenario:} MRI brain scan with $k=4$ tissue types (gray matter, white matter, CSF, tumor). Each pixel must be assigned to one type.
\textbf{Model:} Graph where vertices = pixels, edges = adjacency. Terminals = seed pixels labeled by radiologist for each tissue type. Edge weights = similarity (inverse intensity difference). Multiway cut separates the 4 terminals with minimum boundary cost = maximum coherence within each tissue region.
\textbf{Why multiway cut:} Unlike $k$-means, multiway cut respects boundary continuity — cuts follow intensity gradients. The $2-2/4 = 1.5$-approximation guarantees the segmentation boundary is at most $1.5\times$ the optimal length.
\textbf{In practice:} With pre-seeded terminals from a CNN, the multiway cut refines boundaries better than pure CNN output (post-processing step in medical imaging pipelines).
\textbf{$k$-cut variant:} No pre-labeled seeds — find $k$ components of minimum boundary. Used for unsupervised segmentation.

\section{Tight Example}

$2k$ vertices in a cycle: $k$ terminals alternating with $k$ Steiner vertices. Terminal-Steiner edges weight $2-\epsilon$, Steiner-Steiner edges weight 1. Algorithm picks $k-1$ heavy edges; optimal picks $k$ light edges. Ratio $\to 2-2/k$.

\section{Exercises}
\begin{exercise}
Show Algorithm 4.3 can find a $k$-cut within $2-2/k$ using multiway cut as subroutine.
\end{exercise}
\begin{exercise}
Prove the greedy multiway cut algorithm (iteratively remove lightest $s_i$-$s_j$ cut) also achieves $2-2/k$.
\end{exercise}

